\chapter{Aeroelastic flutter without rational approximations}\label{ch:aero}

Flutter is a dynamic instability of a flexible structure in an airflow:
above a critical speed the aerodynamic forces feed energy into a structural
mode, which then grows. Its analysis is a feedback-stability problem with an
\emph{irrational} feedback path, Theodorsen's function. The standard
practice approximates that function by a rational one, so that the
aeroelastic system becomes a finite state-space model
\citep{edwards1977,kaiser2022}. This chapter shows how to decide stability
and compute the flutter speed with the exact function instead, using the
two-degree-of-freedom typical section and two published benchmarks
\citep{perry2015,kaiser2022}. The code is in
\path{examples/aeroelasticity/} and \path{examples/models/}.

\section{The typical section as a feedback system}

The section (Figure~\ref{fig:section}) has semichord $b$, the elastic axis
at $ab$ from midchord and the centre of gravity at $x_\alpha b$ behind the
elastic axis. With plunge $h$ (positive downwards) and pitch $\alpha$
(positive nose-up) about the elastic axis, $q=[h,\alpha]^T$,
$S_\alpha=mx_\alpha b$ and $I_\alpha=mr_\alpha^2b^2$, the equations per unit
span are
\begin{equation}
  M\ddot q+Kq=\begin{bmatrix}-L\\ M_a\end{bmatrix},\qquad
  M=\begin{bmatrix}m&S_\alpha\\S_\alpha&I_\alpha\end{bmatrix},\qquad
  K=\begin{bmatrix}k_h&0\\0&k_\alpha\end{bmatrix},
\end{equation}
with the lift $L$ (upwards) and the nose-up aerodynamic moment $M_a$ about
the elastic axis.

\begin{figure}[H]
\centering
\begin{tikzpicture}[>=Latex,thick,font=\small,x=1.05cm,y=1.05cm,
  spring/.style={decorate,decoration={zigzag,segment length=4pt,amplitude=3.5pt}}]
  % Airfoil (symmetric, leading edge at x=-3, trailing edge at x=3; chord 2b)
  \fill[blue!10] (-3,0) .. controls (-2.95,0.42) and (-1.9,0.62) .. (-0.6,0.58)
       .. controls (0.8,0.52) and (2.2,0.22) .. (3,0)
       .. controls (2.2,-0.22) and (0.8,-0.52) .. (-0.6,-0.58)
       .. controls (-1.9,-0.62) and (-2.95,-0.42) .. (-3,0) -- cycle;
  \draw[blue!60!black] (-3,0) .. controls (-2.95,0.42) and (-1.9,0.62) .. (-0.6,0.58)
       .. controls (0.8,0.52) and (2.2,0.22) .. (3,0)
       .. controls (2.2,-0.22) and (0.8,-0.52) .. (-0.6,-0.58)
       .. controls (-1.9,-0.62) and (-2.95,-0.42) .. (-3,0) -- cycle;
  \draw[gray,dashed,thin] (-3.3,0)--(3.4,0);
  % Reference points: midchord (0), elastic axis (a b, here a<0), centre of gravity
  \coordinate (EA) at (-1.2,0);
  \coordinate (CG) at (-0.45,0);
  \draw[thin] (0,-0.12)--(0,0.12);
  \node[font=\scriptsize] at (0,-0.85) {midchord};
  \draw[thin,gray] (0,-0.7)--(0,-0.14);
  \fill[black] (CG) circle (2.2pt);
  \node[font=\scriptsize,above] at (-0.45,0.05) {CG};
  \filldraw[fill=white,draw=orange!80!black,very thick] (EA) circle (3.4pt);
  \node[font=\scriptsize,above left] at (-1.25,0.05) {EA};
  % Plunge spring k_h from the elastic axis to the ground
  \draw (EA)--(-1.2,-0.8);
  \draw[spring] (-1.2,-0.8)--(-1.2,-2.0);
  \draw (-1.2,-2.0)--(-1.2,-2.2);
  \draw (-1.7,-2.2)--(-0.7,-2.2);
  \foreach \x in {-1.65,-1.45,...,-0.75} \draw[thin] (\x,-2.2)--(\x-0.15,-2.38);
  \node[right] at (-1.0,-1.4) {$k_h$};
  % Torsional spring k_alpha (spiral around the elastic axis)
  \draw[orange!80!black] (-1.2,-0.32) arc[start angle=-90,end angle=160,radius=0.32]
     arc[start angle=160,end angle=300,radius=0.45];
  \node[orange!80!black,right] at (-0.95,-0.45) {$k_\alpha$};
  % Free stream
  \draw[->,blue!60!black] (-5.1,0.35)--(-3.7,0.35) node[midway,above] {$U$};
  % Degrees of freedom
  \draw[->,red!70!black] (3.7,0.6)--(3.7,-0.6) node[right] {$h$ (down)};
  % Nose-up pitch: the leading edge (left) moves up, i.e. clockwise here.
  \draw[->,red!70!black] (-2.8,-0.62) arc[start angle=201,end angle=148,radius=1.72]
     node[above left] {$\alpha$ (nose-up)};
  % Loads
  \draw[->,accent,very thick] (-1.9,-0.2)--(-1.9,1.2) node[above] {$L$};
  % Dimensions
  \draw[<->,thin] (-3,2.2)--(0,2.2) node[midway,above] {$b$};
  \draw[<->,thin] (0,2.2)--(3,2.2) node[midway,above] {$b$};
  \draw[thin,gray] (-3,2.1)--(-3,0.1) (3,2.1)--(3,0.1) (0,2.1)--(0,0.7);
  \draw[<->,thin] (-1.2,-2.75)--(0,-2.75);
  \node[font=\scriptsize,right] at (0.05,-2.75) {$a\,b$ \ ($a<0$: elastic axis ahead of midchord)};
  \draw[thin,gray] (0,-2.65)--(0,-1.0);
  \draw[<->,thin] (-1.2,1.55)--(-0.45,1.55) node[midway,above,font=\scriptsize] {$x_\alpha b$};
  \draw[thin,gray] (-1.2,1.45)--(-1.2,0.15) (-0.45,1.45)--(-0.45,0.1);
\end{tikzpicture}

\caption{Typical section. The springs $k_h$ and $k_\alpha$ act at the
elastic axis (EA); $L$ is the lift.}
\label{fig:section}
\end{figure}

For a motion $e^{st}$, Theodorsen's theory (incompressible, inviscid,
two-dimensional flow, small motions, Kutta condition) gives
\begin{align}
  L&=\pi\rho b^2\,[s^2h+Us\alpha-abs^2\alpha]+2\pi\rho Ub\,C(p)\,w,\\
  M_a&=\pi\rho b^2\,[abs^2h-Ub(\tfrac12-a)s\alpha-b^2(\tfrac18+a^2)s^2\alpha]
       +2\pi\rho Ub^2(a+\tfrac12)\,C(p)\,w,
\end{align}
where $w=sh+[U+b(\tfrac12-a)s]\alpha$ is the downwash at three-quarter chord
and $p=sb/U$. The first terms are non-circulatory (added mass and damping
of the air); the circulatory terms carry the memory of the wake through
Theodorsen's function $C$. Writing the loads as $A(s,U)q$ gives the dynamic
stiffness matrix and the characteristic function
\begin{equation}
  D(s,U)=s^2M+K-A(s,U),\qquad \Delta(s,U)=\det D(s,U),
  \label{eq:aerodelta}
\end{equation}
whose roots are the aeroelastic modes, in the same way as the eigenvalues of
a state-space matrix are the roots of $\det(sI-A)$. Explicitly,
$A=\pi\rho U^2(p^2A_2+pA_1+A_0)$ with
\begin{gather*}
A_2=\begin{bmatrix}-1&ab\\ab&-(\tfrac18+a^2)b^2\end{bmatrix},\qquad
A_0=\begin{bmatrix}0&-2Cb\\0&2C(\tfrac12+a)b^2\end{bmatrix},\\
A_1=\begin{bmatrix}-2C&[-1-2C(\tfrac12-a)]\,b\\2C(\tfrac12+a)\,b&(\tfrac12-a)[2C(\tfrac12+a)-1]\,b^2\end{bmatrix},
\end{gather*}
in agreement with \citet[Eq.~(33)]{kaiser2022}. At $U=0$ the limit is
$D=s^2(M+M_{\mathrm{air}})+K$ with the still-air added mass
$M_{\mathrm{air}}=\pi\rho b^2\left[\begin{smallmatrix}1&-ab\\-ab&(1/8+a^2)b^2\end{smallmatrix}\right]$.

\section{Theodorsen's function in the Laplace domain}

The generalized Theodorsen function is \citep{edwards1977}
\begin{equation}
  C(p)=\frac{K_1(p)}{K_0(p)+K_1(p)},
\end{equation}
with the modified Bessel functions $K_0$, $K_1$. On the imaginary axis,
$p=\ii k$ with the reduced frequency $k=\omega b/U$, it reduces to the
classical $C(k)=H_1^{(2)}(k)/[H_1^{(2)}(k)+\ii H_0^{(2)}(k)]$; $C(0)=1$ and
$C\to\tfrac12$ at high frequency. The implementation uses exponentially
scaled Bessel functions, whose common factor cancels in the ratio.

$C$ is analytic in the plane cut along the negative real axis; $s=0$ is a
logarithmic branch point, and $C(0)=1$ is a limit, not an analytic value.
Therefore $\Delta(\cdot,U)$ is analytic on any region that does not touch
$(-\infty,0]$, and in particular on
\begin{itemize}
  \item the right half-plane $\Real s>0$, which contains the positive real
        axis --- this window counts the unstable roots, including real ones;
  \item the upper half-plane $\Imag s>0$ --- this window shows the stable
        oscillatory modes; since $\Delta(\bar s)=\overline{\Delta(s)}$ the
        lower half-plane holds the conjugates.
\end{itemize}
The branch cut also produces a non-exponential contribution to transient
responses \citep{edwards1977}; the stability statements below concern the
roots of~\eqref{eq:aerodelta}.

\begin{figure}[H]
\centering
\includegraphics[width=\textwidth]{theodorsen_and_domain.png}
\caption{Left: real and imaginary parts of $C(\ii k)$ and of Jones'
approximation. Centre: the approximation error. Right: the branch cut and
the two valid search regions.}
\label{fig:theodorsen}
\end{figure}

\section{Stability in three steps}

The function $\Delta$ is passed to ContourRoots as a handle, with
\code{AssumeAnalytic} justified by the preceding section. The scaled
determinant $\det(S^{-1}DS^{-1})$, $S=\operatorname{diag}(\sqrt{k_h},\sqrt{k_\alpha})$,
is used for conditioning; functions such as $|\det D|$ or the smallest
singular value are not analytic and must not be used.

\paragraph{1. Stability at a given speed.} The number of roots with
$\Real s>0$, the analogue of \code{any(real(eig(A))>0)}:
\begin{lstlisting}
model = aeroelastic_section('nasa');
Delta = @(s,U) aeroelastic_delta(s,U,model);
[p,info] = croots(@(s) Delta(s,170), [1e-3 60 -250 250], 'AssumeAnalytic', true);
\end{lstlisting}
At $U=170$\,ft/s the count is zero (stable); at $180$\,ft/s it is two, a
pair $2.132\pm74.78\ii$. For large $|s|$ the inertia terms dominate,
$D\approx s^2(M+M_{\mathrm{air}})$, so a window of a few times the
structural frequencies suffices; enlarging it is a cheap check. When a mode
is very close to the imaginary axis --- at very low speed, where the air
adds little damping, or within a fraction of a percent of flutter --- the
count needs more contour samples; it is then reported as
\code{unresolved} until \code{ContourRefinements} is increased.

\paragraph{2. Modes.} The window $[-100,50]\times[0.5,180]$ gives, at
170\,ft/s, the modes $-1.1504+75.8683\ii$ ($\zeta=0.015$) and
$-31.4233+63.4831\ii$ ($\zeta=0.444$).

\paragraph{3. Flutter speed.} The spectral abscissa
$\alpha(U)=\max\Real s$ over the modes crosses zero at the flutter speed,
which \code{fzero} locates:
\begin{lstlisting}
upper = [-100 50 0.5 180];
alpha = @(U) max(real(croots(@(s) Delta(s,U), upper, 'AssumeAnalytic', true)));
Uf = fzero(alpha, [150 200]);            % 173.2624 ft/s
\end{lstlisting}
Alternatively, \code{aeroelastic\_flutter} solves
$\Real\Delta(\ii\omega,U)=\Imag\Delta(\ii\omega,U)=0$ for $(\omega,U)$ by
Newton's method and returns $\Real(ds/dU)=-\Real(\Delta_U/\Delta_s)>0$ at the
crossing.

\paragraph{Divergence.} With $C(0)=1$, static divergence occurs at
$U_d=\sqrt{k_\alpha/[2\pi\rho b^2(a+\frac12)]}$ for $a>-\frac12$, when a real
root crosses into the right half-plane through $s=0$. Real roots are
invisible to the upper window but not to the right-half-plane count: at
$1.05\,U_d$ the count finds the real root $s=1.888$ (NASA case).

\section{Benchmarks}

The NASA standard case \citep[Appendix C]{perry2015} recomputes the
flexure--torsion example of NACA Report 496: $\pi\rho b^2/m=0.1$, $a=-0.4$,
$x_\alpha=0.2$, $r_\alpha^2=0.25$, uncoupled frequencies
$\omega_h=\sqrt{k_h/m}=50$ and $\omega_\alpha=\sqrt{k_\alpha/I_\alpha}=100$\,rad/s,
$b=1$\,ft. The dimensional case of \citet[Table 1]{kaiser2022} is in SI
units. Table~\ref{tab:aerobench} compares the results; the published values
are never used as starting points.

\begin{table}[H]\centering\small
\caption{Flutter of the two benchmarks.}\label{tab:aerobench}
\begin{tabular}{lrrrrr}\toprule
Case & Published $U_f$ & Computed $U_f$ & $\omega_f$ [rad/s] & $k_f$ & $U_d$\\\midrule
NASA & 173.26 ft/s, $k_f=0.4355$ & 173.2624 ft/s & 75.462 & 0.43554 & 353.55 ft/s\\
Kaiser--Quero & 212.2 m/s & 212.1729 m/s & 58.438 & 0.27542 & 394.69 m/s\\
\bottomrule\end{tabular}
\end{table}

\begin{figure}[H]
\centering
\includegraphics[width=\textwidth]{nasa_study.png}
\caption{NASA case. Top: root loci (exact, full lines; p-k, dashed) and
frequencies versus speed. Bottom left: largest real part of the modes for
exact, Jones and quasi-steady aerodynamics. Bottom right: difference in
growth rate between p-k and exact.}
\label{fig:nasastudy}
\end{figure}

\begin{figure}[H]
\centering
\includegraphics[width=\textwidth]{dlr_study.png}
\caption{The dimensional benchmark of \citet{kaiser2022}, same layout as
Figure~\ref{fig:nasastudy}.}
\label{fig:dlrstudy}
\end{figure}

\section{What rational approximations change}

\paragraph{Jones' approximation.} The two-lag approximation
\[
  C_J(p)=1-\frac{0.165\,p}{p+0.0455}-\frac{0.335\,p}{p+0.3}
\]
is rational; with two aerodynamic lag states the system becomes a
six-state model whose eigenvalues are computed with \code{eig}
(\code{aeroelastic\_rational\_roots}). At 170\,ft/s it gives
$-1.0467\pm75.0313\ii$ and $-32.4827\pm66.1864\ii$, plus two real lag roots
($-7.12$ and $-34.96$), against the exact $-1.1504\pm75.8683\ii$ and
$-31.4233\pm63.4831\ii$: the critical mode has about 9\,\% less damping. The
flutter speeds are close (Table~\ref{tab:aeroapprox}); a small error in the
flutter speed does not imply small errors in damping elsewhere.

\paragraph{Quasi-steady aerodynamics.} Setting $C=1$ while keeping the
non-circulatory terms ignores the wake memory. For the NASA case it predicts
flutter at 100.01\,ft/s; the dimensional case is already unstable at
$0.01$\,m/s.

\paragraph{The p-k method.} Evaluating the whole aerodynamic matrix at
$p=\ii\,\Imag(s)b/U$ while keeping the complex $s$ in the structural terms
\citep[Eq.~(4)]{kaiser2022} gives the same equation as the exact method on
the imaginary axis, hence the same flutter speed, but different damping away
from it (Figure~\ref{fig:nasastudy}). The p-k target is not analytic in $s$
and is solved by Newton's method in two real variables, never by contour
counting.

\begin{table}[H]\centering\small
\caption{Flutter speed with approximate aerodynamics.}\label{tab:aeroapprox}
\begin{tabular}{lrrr}\toprule
Case & Exact & Jones (difference) & Quasi-steady\\\midrule
NASA [ft/s] & 173.2624 & 172.8573 ($-0.23$\,\%) & 100.01\\
Kaiser--Quero [m/s] & 212.1729 & 210.9854 ($-0.56$\,\%) & unstable at $0.01$\\
\bottomrule\end{tabular}
\end{table}

\section{Parameter study}

Varying only the plunge stiffness of the NASA case, so that
$\omega_h/\omega_\alpha$ goes from 0.30 to 0.80, lowers the flutter speed from
197.8 to 135.7\,ft/s: flutter is easier as the uncoupled frequencies
approach each other (Figure~\ref{fig:aeroparam}). Every point is confirmed
by contour counts at 98\,\% and 102\,\% of its flutter speed.

\begin{figure}[H]
\centering
\includegraphics[width=0.85\textwidth]{nasa_parameter_study.png}
\caption{Flutter speed and reduced frequency versus the uncoupled frequency
ratio (NASA case).}
\label{fig:aeroparam}
\end{figure}

\section{Validation and limitations}

The study \path{examples/aeroelasticity/run_aeroelastic_study.m} and the
regression tests check: the published flutter values; the Bessel form of $C$
against the Hankel form, its limits and conjugate symmetry; the lift and
moment formulas assembled independently of the matrix implementation; the
still-air added mass; complete contour counts with two modes in the upper
window at every speed of the sweeps, with enlarged windows and doubled
sampling at selected speeds; zero unstable roots below and two above flutter
wherever the count is conclusive, and a real unstable root above divergence;
an independent harmonic flutter formulation (eigenvalues in $U^2$ of
$Kq=U^2B(k)q$ with Hankel functions), which agrees to $10^{-6}$ in speed; and
the Jones state-space eigenvalues against the Jones characteristic function.

The model is linear, incompressible and two-dimensional; ``exact'' refers to
evaluating it without rational approximation. The counts are complete within
the stated rectangles; flutter candidates from Newton's method are confirmed
by the sweeps and counts, not by a global search in the $(\omega,U)$ plane.

\begin{remark}
This example required one change in the solver: Newton steps and
finite-difference stencils are now kept inside the current cell, so that a
function is never evaluated outside the search region. Before, a
right-half-plane search could step onto the branch cut of $C$.
\end{remark}
